\Chapter{CONCLUSION}\label{sec:Conclusion}

\section{Summary of Works}

The objective of this thesis was to improve the robustness of 3D multilayer ice accretion simulations, with particular emphasis on the numerical difficulties associated with droplet impingement, repeated geometry evolution, and the treatment of increasingly complex iced surfaces. The work was conducted within the \textsc{CHAMPS} icing framework and resulted in developments addressing three complementary components of the icing simulation framework: Lagrangian droplet tracking, level-set-based geometry evolution, and immersed-boundary ice accretion.

The first contribution concerned the development and assessment of a Lagrangian particle-tracking methodology for droplet impingement. A face-to-face trajectory algorithm was combined with ray-tracing-enhanced intersection detection, adaptive seeding, and surface interpolation strategies to improve the robustness and computational efficiency of the method. Supercooled Large Droplet effects were incorporated into the droplet formulations, including models accounting for droplet deformation, splashing, and secondary-droplet transport. The resulting predictions showed good agreement with reference and experimental data. Comparisons between the Eulerian and Lagrangian formulations further demonstrated that both approaches converge toward consistent collection-efficiency and collected-water-mass predictions.

The comparison nevertheless highlighted important differences between the two formulations. The Eulerian approach naturally follows the decomposition of the computational mesh and provides a predictable workload and efficient parallelization for large 3D simulations. The Lagrangian formulation can be considerably less expensive when a relatively small number of trajectories is sufficient, particularly in 2D calculations, but becomes more challenging on large distributed-memory configurations because of irregular trajectory lengths, load balancing, and access to mesh information distributed across computational domains. Its principal advantage lies in retaining the history and impact location of individual droplets, providing a natural framework for more detailed physical models involving splashing and re-impingement of the resulting droplets. The two formulations should therefore be regarded as complementary rather than competing approaches.

The central contribution of this thesis was the development of a continuous level-set methodology for 3D multilayer ice accretion. The evolving ice--air interface is represented implicitly as the zero contour of a signed-distance function defined on an auxiliary volume mesh. The proposed methodology combines surface-quantity propagation, semi-implicit level-set advection, local interface reconstruction for ice-volume evaluation, interface extraction, and surface smoothing. The method was assessed using representative icing benchmarks and demonstrated the ability to reproduce complex ice geometries while maintaining a conformal reconstructed surface.

The comparison with explicit algebraic surface displacement illustrates the principal motivation for the level-set approach. Direct displacement remains simple and computationally attractive for single-layer calculations and geometries undergoing limited deformation. However, its robustness progressively depends on the quality and connectivity of the evolving surface mesh as additional accretion layers are introduced. The implicit level-set representation removes this direct dependency during geometry evolution and consequently allows more complex geometries to develop without surface-element entanglement. Although this requires additional operations associated with volume-meshing, field propagation, advection, and interface extraction, these procedures can largely be automated and accelerated. The additional numerical complexity therefore provides an advantageous trade-off when robust multilayer evolution is required.

The third contribution extended the icing framework to an immersed-boundary representation. The methodology combines immersed-boundary aerodynamic and droplet treatments with the existing thermodynamic model and level-set geometry-evolution procedure. Verification on canonical configurations demonstrated consistency between body-fitted and immersed-boundary droplet predictions, while rime and glaze ice-accretion calculations showed encouraging agreement between the two approaches. These results demonstrate the feasibility of performing the complete icing sequence without requiring the generation of a new aerodynamic volume mesh at each accretion layer.

The immersed-boundary investigation also showed that eliminating body-fitted remeshing does not eliminate the mesh-resolution problem. Instead, the numerical challenge shifts toward maintaining sufficient spatial resolution around the evolving immersed interface. Prescribing a sufficiently large region of uniformly fine cells can rapidly become computationally expensive for 3D configurations. Adaptive mesh refinement therefore represents a natural complement to the immersed-boundary approach, allowing the computational resolution to follow the evolving ice geometry without requiring excessive refinement throughout the complete domain.

Taken together, the developments presented in this thesis address different but interconnected limitations of 3D icing simulations. Lagrangian particle tracking increases the flexibility available for modeling detailed droplet physics, the level-set formulation improves the robustness of multilayer geometry evolution, and the immersed-boundary methodology reduces the dependency of the flow solution on repeated geometry-conforming volume remeshing. Among these developments, the level-set formulation provides the central geometry-evolution framework around which the other methodologies can be combined. The resulting numerical framework provides a foundation for increasingly automated and robust 3D multilayer ice accretion simulations.

\section{Limitations}\label{sec:Limitations}

Despite the developments achieved in this work, several limitations remain within the different components of the icing framework. The accuracy of ice-accretion predictions remains dependent on the physical models preceding geometry evolution. A robust geometric methodology cannot compensate for inaccuracies in aerodynamic quantities, droplet impingement, heat transfer, surface roughness, or thermodynamic modeling. This limitation is particularly important under glaze ice conditions, for which horn development exhibits substantial sensitivity to the convective heat-transfer prediction and the prescribed equivalent sand-grain roughness.

The Lagrangian formulation also retains computational limitations for large 3D configurations. The number of trajectories required to obtain a converged impingement field is not known a priori, and individual trajectories can have substantially different computational costs. Adaptive seeding reduces the number of unnecessary trajectories, but efficient execution on distributed-memory systems continues to depend on appropriate load balancing and data-locality management. Consequently, the computational advantage observed for smaller or shared-memory problems does not necessarily extend directly to large distributed-memory calculations.

% The level-set methodology achieves greater geometric robustness at the expense of an auxiliary volume mesh and additional field-initialization, propagation, advection, interface-extraction, and smoothing operations. The accuracy of the reconstructed geometry consequently remains dependent on the resolution of the auxiliary mesh and on the interface-processing procedures. Furthermore, quantitative validation of highly 3D ice morphologies remains challenging because available experimental measurements are predominantly based on 2D cuts or maximum-combined-cross-section representations. The generation of complex 3D structures therefore does not necessarily imply that their detailed spanwise morphology is physically accurate.

Finally, the present immersed-boundary implementation remains limited by the absence of dynamic adaptive mesh refinement. Adequate resolution must currently be prescribed over a sufficiently large region to accommodate the expected evolution of the ice geometry. This requires prior knowledge or an estimate of the resulting ice shape and can lead to unnecessarily large computational meshes when refinement is maintained in regions that are not ultimately reached by the evolving interface. The development of an efficient adaptive strategy capable of refining and coarsening the background mesh as the ice surface evolves therefore represents an important step toward improving the computational efficiency of the immersed-boundary approach for large-scale 3D multilayer icing simulations.

\section{Future Research}

The developments presented in this thesis suggest several directions for future research. For Lagrangian droplet tracking, future developments should focus on improving computational efficiency for large 3D configurations, particularly through GPU acceleration. A GPU-oriented implementation would, however, require data structures and algorithms specifically designed to accommodate irregular memory-access patterns, mesh-intersection searches, and heterogeneous trajectory lengths. Combined with the adaptive cloud-refinement and load-balancing strategies developed in this work, such an implementation could substantially reduce the computational cost of large-scale 3D impingement calculations.

Further development of the level-set methodology should focus on improving the geometry-evolution workflow for large 3D icing simulations. MMG is a general-purpose meshing tool and is not specifically designed to generate surface meshes for CFD simulations. A dedicated surface-meshing strategy could provide greater control over surface-element size and quality, particularly in highly curved regions of the evolving ice geometry, while ensuring that the extracted surface remains suitable for subsequent volume-mesh generation and flow calculations. The current methodology also enforces conservation of the total ice volume at the global level. Future work should investigate local ice-mass conservation strategies to ensure that the spatial distribution of the accreted mass is preserved during level-set advection and interface reconstruction. More generally, the increased robustness of the level-set geometry-evolution framework provides an opportunity to investigate physical modeling assumptions more systematically. In particular, the influence of spatially and temporally varying surface roughness, variable ice-density formulations, and their interaction with the thermodynamic model can now be assessed with greater confidence over complex 3D multilayer configurations, without the robustness of the geometry-update procedure becoming a limiting factor.

Finally, further development of the immersed-boundary icing framework should focus on improving the near-wall heat-transfer treatment and incorporating adaptive mesh refinement. Improving the near-wall heat-transfer treatment is necessary to further bridge the gap between immersed-boundary and established body-fitted approaches for icing simulations. Adaptive mesh refinement would dynamically maintain the spatial resolution required near the immersed interface as the ice geometry evolves, while retaining coarser cells elsewhere. Together, these developments would improve the applicability of the immersed-boundary approach within \textsc{CHAMPS} for complex 3D multilayer icing simulations.